\documentclass[main.tex]{subfiles}
\begin{document}
\chapter{Jet Bundles and Higher-Order Obstructions} \label{chapter:jets}

%% All commands are in macros.sty

\subsection{Jet Obstructions}


Here we describe one approach to obstructing unirationality. The idea is very simple: recall that in characteristic zero we may use forms to obstruct unirationality. However, the issue is that degree $p$ purely inseparable maps can kill first-order derivatives. The fundamental observation is that they cannot kill $p^{\text{th}}$-order derivatives. Therefore, if we can find higher-order forms on our varieties then these will pull back nontrivially even along certain inseparable maps to produce obstructions. This idea can be formalized using various notions of jet bundles. We plan to extend and combine these techniques with the foliation--descent methods developed in this thesis to develop more applicable higher-jet unirationality obstructions. 

\begin{defn}
Let $X \to S$ be a smooth scheme. Then $J^n_{X/S} := \pi_{1*} \pi_2^* \struct{X}$ where the projections are along the $n^{\text{th}}$-order thickening of the diagonal $\Delta^n_{X/S} := V(\I_{\Delta}^{n+1}) \subset X \times_S X$. This satisfies the following universal property,
\[ \shHom_{\struct{X}}(J^n_{X/S},\E) = \shDiff^{\le n}_{X}(\struct{X}, \E) \]
where $\shDiff^{\le n}_{X}(\struct{X}, \E)$ is the sheaf of differential operators $D : \cO_X \to \E$ of order $\le n$ (see \cite[IV$_4$, \S~16.8]{EGAIV}).
\end{defn}

We begin by defining a notion of bigness for lifts of jets.

\begin{defn}
Let $\wt{J}^n_X := \ker{(J^n_X \to \struct{X})}$ using the natural projection $J^n_X \to J^0_X = \struct{X}$.
\end{defn}

\begin{defn}
Let $X$ be a smooth proper variety. Define the \textit{jet-amplitude} of $X$ to be
\[ \jetamp(X) := \sup \left\{ \frac{n}{r} : \exists \, \omega \in H^0(X, \wt{J}^{n}_X) \text{ such that } \omega|_{\wt{J}^{r}_X} \neq 0 \right\} \]
\end{defn}

\begin{theorem}[C]
Let $X$ be a smooth proper variety over a perfect field $k$. Suppose that $\log_p \jetamp(X) > \ell$ then there does not exist a unirational parametrization $\P^n \rat X$ of inseparability degree $\ell$.
\end{theorem}

\begin{proof}
There is a factorization of rational maps
\begin{center}
\begin{tikzcd}
X^{(1/q)} \arrow[r, dashed, "f'"] \arrow[rd, "F"'] & \P^n \arrow[d, dashed, "f"]
\\
& X
\end{tikzcd}
\end{center}
where $F : X^{(1/q)} \to X$ is the $\ell^{\text{th}}$-relative Frobenius.
Using extension in codimension $\ge 2$, there are well-defined pullback maps
\begin{center}
\begin{tikzcd}
H^0(X, \wt{J}^{s}_X) \arrow[r] \arrow[d, "f^*"] & H^0(X, \wt{J}^{r}_X) \arrow[d, "f^*"]
\\
H^0(\P^n, \wt{J}^{s}_{\P^n}) \arrow[d, "f'^*"] \arrow[r] & H^0(\P^n, \wt{J}^{r}_{\P^n}) \arrow[d, "f'^*"]
\\
H^0(X^{(1/q)}, \wt{J}^{s}_{X^{(1/q)}}) \arrow[r] & H^0(X^{(1/q)}, \wt{J}^{r}_{X^{(1/q)}})
\end{tikzcd}
\end{center}
where the composition along the column is $F^*$. Therefore, it suffices to show that $F^* \omega \neq 0$. To do this, we may work \etale-locally and reduce to the following computation: $J^r_{\A^n}$ is represented by the algebra,
\[ k[x_1, \dots, x_n][\d{x_1}, \dots, \d{x_n}] / (\d{x_1}, \dots, \d{x_n})^{r+1} \] and $F^* \d{x_i} = \d{x_i^{q}} = (\d{x_i})^{q}$. Therefore, this is nonzero if $s/r \ge q$.
\end{proof}

Unfortunately it is quite difficult to compute $\jetamp(X)$. For ``generic'' complete intersections of large degree, it is known that $\jetamp(X) = \infty$. However, for any particular example it is very hard to get one's hands on this thing. There is a trivial bound as follows.


\begin{prop}
Let $X$ be a smooth projective variety. Then,
\[ \jetamp(X) \ge \left(\frac{\ul{\hat{h}}^0(\Omega_X)}{\hat{h}^1(\Omega_X)}\right)^{\frac{1}{2 \dim{X}}} \]
where $\hat{h}^i$ are the so-called \textit{asymptotic cohomology}
\[ \hat{h}^i(X, \E) := \limsup_{m \to \infty} \frac{h^0(X, \nSym{m}{\E})}{m^{\dim{X} + \rank{\E} - 1}} \]
and likewise $\ul{\hat{h}}^i$ are the \textit{lower asymptotic cohomology}
\[ \ul{\hat{h}}^i(X, \E) := \liminf_{m \to \infty} \frac{h^0(X, \nSym{m}{\E})}{m^{\dim{X} + \rank{\E} - 1}} \]
\end{prop}


\begin{proof}
This arises from the exact sequence
\begin{center}
\begin{tikzcd}
0 \arrow[r] & \nSym{n}{\Omega_X} \arrow[r] & J^n_X \arrow[r] & J^{n-1}_X \arrow[r] & 0
\end{tikzcd}
\end{center}
\end{proof}

\begin{example}
Suppose $h^0(X, \nSym{2}{\Omega_X}) \ge h^1(X, \nSym{3}{\Omega_X}) + h^1(X, \nSym{4}{\Omega_X})$ then $X$ is not a Zariski space in characteristic $2$.
\end{example}


\subsection{Proof of the Jet Amplitude Theorem}

\begin{theorem}[C]
Let $X$ be a smooth proper variety over a perfect field $k$. Suppose that $\log_p \jetamp(X) > \ell$ then there does not exist a unirational parametrization $\P^n \rat X$ of inseparability degree $\ell$.
\end{theorem}

\begin{proof}
Since $\jetamp$ only increases along separable maps we may assume that $X$ has a purely inseparable unirational parametrization $f : \P^n \rat X$ of degree $k$. Now suppose that $\jetamp \ge p^{\ell} := q$ then there are integers $r,s \ge 0$ such that $\omega \in H^0(X, \wt{J}_X^{s})$ such that $\omega|_{\wt{J}^{r}} \neq 0$ and $s \ge q r$. By decreasing $r$ we may assume that $\omega|_{\wt{J}^{r-1}} = 0$ so $\omega|_{\wt{J}^{r}} \in H^0(X, \nSym{r}{\Omega_X})$ then decreasing $s$ we may assume $s = q r$. There is a factorization of rational maps
\begin{center}
\begin{tikzcd}
X^{(1/q)} \arrow[r, dashed, "f'"] \arrow[rd, "F"'] & \P^n \arrow[d, dashed, "f"]
\\
& X
\end{tikzcd}
\end{center}
where $F : X^{(1/q)} \to X$ is the $\ell^{\text{th}}$-relative Frobenius.
Using extension in codimension $\ge 2$, there are well-defined pullback maps
\begin{center}
\begin{tikzcd}
H^0(X, \wt{J}^{s}_X) \arrow[r] \arrow[d, "f^*"] & H^0(X, \wt{J}^{r}_X) \arrow[d, "f^*"]
\\
H^0(\P^n, \wt{J}^{s}_{\P^n}) \arrow[d, "f'^*"] \arrow[r] & H^0(\P^n, \wt{J}^{r}_{\P^n}) \arrow[d, "f'^*"]
\\
H^0(X^{(1/q)}, \wt{J}^{s}_{X^{(1/q)}}) \arrow[r] & H^0(X^{(1/q)}, \wt{J}^{r}_{X^{(1/q)}})
\end{tikzcd}
\end{center}
where the composition along the column is $F^*$. Therefore, it suffices to show that $F^* \omega \neq 0$. To do this, we may work \etale-locally. Indeed, if $c : U \to X$ is an \etale chart then it suffices to show that $F_U^* \omega|_U \neq 0$ since the diagram
\begin{center}
\begin{tikzcd}
U^{(1/q)} \arrow[r, "F_U"] \arrow[d] & U \arrow[d]
\\
X^{(1/q)} \arrow[r, "F"] & X
\end{tikzcd}
\end{center}
commutes. Since $X$ is smooth, there are \etale charts $U \to X$ with $g : U \to \A^n_k$ \etale. The diagram,
\begin{center}
\begin{tikzcd}
\wt{J}^s_{\A^n} \arrow[d, "F^*"] \arrow[r] & \wt{J}^r_{\A^n} \arrow[d, "F^*"] \arrow[ld, "t", dashed]
\\
\wt{J}^s_{\A^n} \arrow[r] & \wt{J}^r_{\A^n}
\end{tikzcd}
\end{center}
has a commuting lift $t$. Indeed, $J^r_{\A^n}$ is represented by the algebra,
\[ k[x_1, \dots, x_n][\d{x_1}, \dots, \d{x_n}] / (\d{x_1}, \dots, \d{x_n})^{r+1} \] and $F^* \d{x_i} = \d{x_i^{q}} = (\d{x_i})^{q}$. Therefore, the kernel of $\wt{J}^s_{\A^n} \to \wt{J}^r_{\A^n}$ is the space of polynomials $h(\d{x_1}, \dots, \d{x_n})$ with coefficients in $k[x_1, \dots, x_n]$ of minimal degree $\ge r + 1$ and hence $F^* h = h(\d{x_1^{q}}, \dots, \d{x_n^{q}}) = 0$ in $J^s_{\A^n}$ since it has minimal degree $\ge q (r + 1) \ge s + 1$. Furthermore, consider a symmetric form $\eta \in H^0(\A^n, \nSym{r}{\Omega})$ write
\[ \eta = \sum_{I} f_I(x_1, \dots, x_n) \d{x_{i_1}} \cdots \d{x_{i_r}} \]
then we see
\[ t(\eta) = \sum_I f_I(x_1^q, \dots, x_n^q) \d{x_{i_1}^q} \cdots \d{x_{i_r}^q} \in H^0(\A^n, \nSym{s}{\Omega}) \]
but for different $I \subset \{1, \dots, n\}$ of size $r$, these are distinct basis elements of $H^0(\A^n, \nSym{s}{\Omega})$. Therefore, $t(\eta) = 0$ if and only if all $f_I(x_1^q, \dots, x_n^q) = 0$ if and only if all $f_I = 0$. Hence $t$ is injective. Since $g$ is \etale, we get isomorphisms $g^* \wt{J}_{\A^n}^r \iso \wt{J}^r_{U}$ for all $r$ so the diagram pulls back to
\begin{center}
\begin{tikzcd}
\wt{J}^s_{U} \arrow[r] \arrow[d, "F^*_U"'] & \wt{J}^r_{U} \arrow[d, "F^*_U"] \arrow[dl, "t_U"']
\\
\wt{J}^s_{U} \arrow[r] & \wt{J}^r_{U}
\end{tikzcd}
\end{center}
such that the restriction $t_U : \nSym{r}{\Omega_U} \to \nSym{s}{\Omega_U}$ is injective. Therefore $F^*_U \omega|_U = t(\omega|_{\wt{J}^r_U}) \neq 0$ because we assumed that $\omega|_{\wt{J}^r} \in H^0(X, \nSym{r}{\Omega_X})$ and is nonzero.
\end{proof}

\begin{prop}
Let $X$ be a smooth projective variety. Then,
\[ \jetamp(X) \ge \left(\frac{\ul{\hat{h}}^0(\Omega_X)}{\hat{h}^1(\Omega_X)}\right)^{\frac{1}{2 \dim{X}}} \]
where $\hat{h}^i$ are the so-called \textit{asymptotic cohomology}
\[ \hat{h}^i(X, \E) := \limsup_{m \to \infty} \frac{h^0(X, \nSym{m}{\E})}{m^{\dim{X} + \rank{\E} - 1}} \]
and likewise $\ul{\hat{h}}^i$ are the \textit{lower asymptotic cohomology}
\[ \ul{\hat{h}}^i(X, \E) := \liminf_{m \to \infty} \frac{h^0(X, \nSym{m}{\E})}{m^{\dim{X} + \rank{\E} - 1}} \]
\end{prop}


\begin{proof}
This arises from the exact sequence
\begin{center}
\begin{tikzcd}
0 \arrow[r] & \nSym{n}{\Omega_X} \arrow[r] & J^n_X \arrow[r] & J^{n-1}_X \arrow[r] & 0
\end{tikzcd}
\end{center}
Therefore,
\[ \dim \im{(H^0(X, J^{rs}_X) \to H^0(X, J^r_X))} \ge h^0(X, J^r_X) - \sum_{k = r+1}^{rs} \left[h^1(X, \nSym{k}{\Omega_X}) \right] \]
but also
\[ h^0(X, J^r_X) \ge \sum_{i = 1}^r \left[ h^0(X, \nSym{i}{\Omega_X}) - h^1(X, \nSym{i}{\Omega_X}) \right] \]
therefore
\[ \dim \im{(H^0(X, J^{rs}_X) \to H^0(X, J^r_X))} \ge \sum_{i = 1}^r h^0(X, \nSym{i}{\Omega_X}) - \sum_{i = 1}^{rs} h^1(X, \nSym{i}{\Omega_X}) \]
Applying the lemma, for any $\epsilon > 0$ for $r \gg_{\epsilon} 0$ we have
\[ \sum_{i = 1}^r h^0(X, \nSym{i}{\Omega_X}) - \sum_{i = 1}^{rs} h^1(X, \nSym{i}{\Omega_X}) \ge \frac{1}{2 \dim{X}} \left[ (\ul{\hat{h}}^0 - \epsilon) r^{2 \dim{X}} - (\hat{h}^1 + \epsilon) (rs)^{2 \dim{X}} \right] \]
Thus,
\[ s^{2 \dim{X}} \le \frac{\ul{\hat{h}}^0 - \epsilon}{\hat{h}^1 + \epsilon} \implies \jetamp(X) \ge s  \]
\end{proof}

\begin{lemma}
Let $f : \N \to \RR^+$ be an integer function. Suppose that
\[ \liminf_{n \to \infty} \frac{f(n)}{n^k} = a \quad \text{ and } \quad \limsup_{n \to \infty} \frac{f(n)}{n^k} = b \]
then\footnote{Notice that the outer inequalities are not usually equalities. For example let
\[ f(n) = \begin{cases} 0 & n \text{ even} \\ 1 & n \text{ odd} \end{cases} \]
then the hypotheses are satisfied with $k = 0$ and $a = 0$ and $b = 1$ but
\[ \sum_{i = 1}^n f(i) = \ceil{\frac{n}{2}} \]
and therefore the limit
\[ \lim_{n \to \infty} \frac{1}{n} \sum_{i = 1}^n f(i) = \frac{1}{2} \]
exists so the inequalities in the conclusion of the lemma are strict. }
\[ \frac{b}{k+1} \ge \limsup_{n \to \infty} \frac{1}{n^{k+1}} \sum_{i = 1}^n f(i) \ge \liminf_{n \to \infty} \frac{1}{n^{k+1}} \sum_{i = 1}^n f(i) \ge \frac{a}{k+1}  \]
\end{lemma}

\begin{proof}
The assumptions give for all $\epsilon > 0$ there is $n_\epsilon$ such that $n \ge n_\epsilon$ implies
\[ (b + \epsilon) n^k \ge f(n) \ge (a - \epsilon) n^k \]
summing this over, we get,
\[ (b + \epsilon) \frac{1}{n^{k+1}} \sum_{i = 1}^n i^k \ge \frac{1}{n^{k+1}} \sum_{i = n_\epsilon}^n f(i) + \frac{1}{n^{k+1}} \sum_{i = 1}^{n_\epsilon} i^k \ge (a - \epsilon) \frac{1}{n^{k+1}} \sum_{i = 1}^n i^k \]
Therefore,
\[ \frac{b + \epsilon}{k + 1} (1 + O(n^{-1})) \ge \frac{1}{n^{k+1}} \sum_{i = 1}^n f(i) + \frac{1}{n^{k+1}} \sum_{i = 1}^{n_\epsilon} (i^k - f(i)) \ge \frac{a - \epsilon}{k+1} (1 + O(n^{-1}))  \]
Taking the limit $n \to \infty$ and then $\epsilon \to 0$ gives the bounds
\[ \frac{b}{k+1} \ge \limsup_{n \to \infty} \frac{1}{n^{k+1}} \sum_{i = 1}^n f(i) \ge \liminf_{n \to \infty} \frac{1}{n^{k+1}} \sum_{i = 1}^n f(i) \ge \frac{a}{k+1}  \]
Furthermore, for any $\epsilon > 0$ and $n_0$ there is $n \ge n_0$ such that $(a + \epsilon) n^k \ge f(n)$
\end{proof}


\ifSubfilesClassLoaded{%
  \bibliographystyle{alpha}
  \bibliography{refs}
}{}

\end{document}
